% Fiche 1 — Comprendre la preuve de la Knowledge Locality (SLD resolution)
% Compilation : pdflatex fiche-preuve.tex   (a besoin de style-fiche.tex)
\documentclass[10pt]{article}
% =====================================================================
%  Style commun aux fiches (bleu marine + bleu pastel)
%  Pour changer les couleurs : modifier uniquement ce fichier.
% =====================================================================
\usepackage[utf8]{inputenc}
\usepackage[T1]{fontenc}
\usepackage[scaled]{helvet}
\renewcommand{\familydefault}{\sfdefault}
\usepackage[a4paper,margin=8mm]{geometry}
\usepackage{xcolor}
\usepackage[most]{tcolorbox}
\usepackage{tikz}
\usepackage{pifont}
\usepackage{enumitem}
\usepackage{array}
\usepackage{etoolbox}
\usepackage{needspace}
\usepackage{amsmath,amssymb}
\usepackage{multirow}
\pagestyle{empty}
\setlength{\parindent}{0pt}

% ---------- Couleurs ----------
\definecolor{marine}{HTML}{3F72B8}   % bleu (plus clair que le marine)
\definecolor{pastel}{HTML}{BFD6F6}   % bleu pastel
\definecolor{fond}{HTML}{F4F7FB}
\definecolor{encre}{HTML}{1F1F1F}
\definecolor{vert}{HTML}{22A55B}
\definecolor{rouge}{HTML}{DC2626}
\pagecolor{fond}
\color{encre}

% ---------- Styles page de synthèse ----------
\newtcolorbox{bloc}[2][]{enhanced, colback=white, colframe=pastel,
  boxrule=0.9pt, arc=0pt, left=5pt, right=5pt, top=3pt, bottom=3pt,
  fontupper=\footnotesize\raggedright, #1,
  before upper={{\large\bfseries\textcolor{marine}{#2}}\par\vspace{2pt}}}
\newtcolorbox{astuce}{colback=pastel!75, colframe=pastel!75, arc=0pt,
  boxrule=0pt, left=4pt, right=4pt, top=2pt, bottom=2pt,
  fontupper=\footnotesize\raggedright}
\newcommand{\num}[1]{\textcolor{marine}{\textbf{#1.}}~}
\newcommand{\fl}{\tikz[baseline=-0.6ex]\draw[->,thick](0,0)--(0.4,0);}
\newcommand{\sous}[1]{\textcolor{marine}{\textbf{#1}}\par}
\setlist[itemize]{leftmargin=9pt, itemsep=0pt, topsep=1pt, parsep=0pt,
  label=\textbullet, before=\raggedright}
\newcolumntype{L}[1]{>{\raggedright\arraybackslash}p{#1}}

% Terme anglais (en gras) avec sa traduction française
\newcommand{\EN}[2]{\textbf{#1}\,{\scriptsize\textcolor{marine}{(#2)}}}

% ---------- Styles pages détaillées ----------
\newtcolorbox{dire}{enhanced, breakable, colback=pastel!25, colframe=marine,
  boxrule=0pt, leftrule=3pt, arc=0pt, left=8pt, right=6pt, top=4pt, bottom=4pt,
  fontupper=\small\itshape, before upper={{\upshape\bfseries\textcolor{marine}{FR}}\enspace}}
\newtcolorbox{direEN}{enhanced, breakable, colback=vert!6, colframe=vert,
  boxrule=0pt, leftrule=3pt, arc=0pt, left=8pt, right=6pt, top=4pt, bottom=4pt,
  fontupper=\small\itshape, before upper={{\upshape\bfseries\textcolor{vert}{EN}}\enspace}}
\newtcolorbox{pourquoi}{colback=fond, colframe=fond, arc=0pt, boxrule=0pt,
  left=6pt, right=6pt, top=3pt, bottom=3pt, fontupper=\small}
\newcommand{\titrebox}[2]{\needspace{10\baselineskip}\vspace{8pt}
  \begin{tcolorbox}[colback=marine, colframe=marine, arc=0pt, boxrule=0pt,
    top=4pt, bottom=4pt, left=8pt]
  \color{white}\Large\bfseries #1.~#2
  \end{tcolorbox}\vspace{2pt}}
\newcommand{\ce}[1]{\begin{pourquoi}\textbf{\textcolor{marine}{À retenir :}} #1\end{pourquoi}}
\newcommand{\st}[1]{\par\needspace{7\baselineskip}\vspace{5pt}{\normalsize\bfseries\textcolor{marine}{#1}}\par\vspace{2pt}}
\newcommand{\sst}[1]{\par\needspace{5\baselineskip}\vspace{3pt}{\small\bfseries\textcolor{encre}{#1}}\par\vspace{1pt}}
\newtcolorbox{moi}{colback=pastel!35, colframe=pastel!35, arc=0pt, boxrule=0pt,
  left=6pt, right=6pt, top=3pt, bottom=3pt, fontupper=\small}
\newtcolorbox{piege}{colback=rouge!6, colframe=rouge, boxrule=0pt, leftrule=3pt, arc=0pt,
  left=8pt, right=6pt, top=3pt, bottom=3pt, fontupper=\small}
% Boîte "énoncé" pour lemmes, propositions, théorèmes
\newtcolorbox{enonce}[1]{enhanced, breakable, colback=white, colframe=marine,
  boxrule=0.8pt, arc=0pt, left=8pt, right=8pt, top=4pt, bottom=4pt,
  fontupper=\small, title={#1}, coltitle=white, colbacktitle=marine,
  fonttitle=\bfseries\small}
% Passage à la mise en page détaillée
\newcommand{\detailpages}{%
  \newpage
  \newgeometry{a4paper,margin=15mm}
  \setlist[itemize]{leftmargin=12pt, itemsep=1pt, topsep=2pt, parsep=0pt,
    label=\textcolor{marine}{\textbullet}}
  \small}


\begin{document}

% ================= EN-TÊTE =================
\begin{minipage}[c]{0.64\textwidth}
  {\fontsize{26}{28}\selectfont\bfseries\textcolor{marine}{PROOF}~CHEAT SHEET 1}\\[1pt]
  {\small Knowledge Locality --- pourquoi un client avec moins de faits tire la même conclusion}
\end{minipage}\hfill
\begin{minipage}[c]{0.34\textwidth}\raggedright
  \normalsize \textbf{Claim} $\to$ \textbf{Tool} (SLD)\\ $\to$ \textbf{Invariant} $\to$ \textbf{Conclusion}
\end{minipage}

\vspace{5pt}

% ================= LIGNE 1 =================
\begin{minipage}[t]{0.32\textwidth}
\begin{bloc}{\num{1}The Lemma (le lemme)}
\sous{Statement (énoncé)}
Pour tout exemple $e$ du client $i$ et tout $H$ :\\[2pt]
\centerline{$\mathcal{B}\cup H\models e \iff \mathcal{B}_i\cup H\models e$}
\sous{Two directions}
\begin{itemize}
  \item $\Leftarrow$ \textbf{easy} : $\mathcal{B}_i\subseteq\mathcal{B}$, ajouter des faits ne retire aucune conclusion (\emph{monotonicity})
  \item $\Rightarrow$ \textbf{hard} : retirer des faits peut faire perdre une preuve. Il faut montrer que la preuve de $e$ n'utilise \textbf{aucun fait manquant}
\end{itemize}
\begin{astuce}
Le sens difficile : celui où l'on \textbf{enlève} des faits.
\end{astuce}
\end{bloc}
\end{minipage}\hfill
\begin{minipage}[t]{0.32\textwidth}
\begin{bloc}{\num{2}The 3 assumptions}
\begin{itemize}
  \item \textbf{Example Locality} : chaque exemple est chez un seul client, et sa constante est locale \fl démarre la preuve chez lui
  \item \textbf{Separated Local Constants} : constantes locales disjointes, aucune dans $\mathcal{R}$ ni $H$ \fl un fait se reconnaît à ses constantes
  \item \textbf{Directions} : chaque argument \texttt{in} est connu avant l'appel \fl on ne sort jamais de l'exemple
\end{itemize}
\begin{astuce}
Chacune sert à un endroit précis de la preuve (box 5).
\end{astuce}
\end{bloc}
\end{minipage}\hfill
\begin{minipage}[t]{0.32\textwidth}
\begin{bloc}{\num{3}SLD in one minute}
\begin{itemize}
  \item \textbf{Goal} $\leftarrow e$ : ``$e$ est faux''. On cherche une contradiction.
  \item \textbf{One step} : choisir un littéral, trouver une clause dont la tête s'unifie, le remplacer par le corps.
  \item \textbf{Success} : but vide = \textbf{refutation} = preuve de $e$.
  \item \textbf{Prolog} = SLD + littéral le plus à gauche + profondeur d'abord + backtracking.
\end{itemize}
\textbf{Sound} : réfutation $\Rightarrow$ conséquence. \textbf{Complete} : conséquence $\Rightarrow$ réfutation (pour \emph{toute} selection rule).
\end{bloc}
\end{minipage}

\vspace{4pt}

% ================= LIGNE 2 =================
\begin{minipage}[t]{0.5\textwidth}
\begin{bloc}{\num{4}SLD on \texttt{f(t800)}}
Clause : \texttt{f(A) :- has\_car(A,B), short(B).}\par\vspace{2pt}
{\scriptsize
\renewcommand{\arraystretch}{1.1}
\begin{tabular}{@{}l L{3.1cm} L{3.9cm}@{}}
\textbf{Step} & \textbf{Goal} & \textbf{Action}\\\hline
0 & \texttt{f(t800)} & clause, \texttt{A=t800}\\
1 & \texttt{has\_car(t800,B), short(B)} & littéral de gauche : fait \texttt{has\_car(t800,t800\_c1)}, \texttt{B=t800\_c1}\\
2 & \texttt{short(t800\_c1)} & aucun fait $\to$ échec, backtracking\\
3 & \texttt{short(t800\_c2)}, \texttt{c3} & même échec\\
4 & \texttt{short(t800\_c4)} & fait présent $\to$ but \textbf{vide}, preuve\\
\end{tabular}}
\vspace{2pt}
\begin{astuce}
\texttt{B} sort de \texttt{has\_car} \emph{avant} d'entrer dans \texttt{short} : c'est la \textbf{direction}.
\end{astuce}
\end{bloc}
\end{minipage}\hfill
\begin{minipage}[t]{0.48\textwidth}
\begin{bloc}{\num{5}The proof in 5 steps}
\begin{enumerate}[leftmargin=12pt, itemsep=1pt, topsep=1pt]
  \item $e$ conséquence $\Rightarrow$ (\textbf{completeness}) il existe une réfutation SLD, littéral de gauche, corps écrits dans l'ordre des directions.
  \item \textbf{Allowed constant for $i$} : locale à $i$, ou écrite dans $\mathcal{R}$ ou $H$. Aucune n'est locale à un autre client (\emph{Separated Local Constants}).
  \item \textbf{Invariant} : à chaque littéral sélectionné, ses arguments \texttt{in} sont des constantes autorisées pour $i$.
  \item Un fait utilisé est \textbf{chez $i$} (sinon un autre client aurait une constante autorisée pour $i$). Une règle utilisée est dans $\mathcal{R}\cup H$.
  \item Même réfutation avec $\mathcal{B}_i\cup H$ $\Rightarrow$ (\textbf{soundness}) $\mathcal{B}_i\cup H\models e$.
\end{enumerate}
\end{bloc}
\end{minipage}

\vspace{4pt}

% ================= LIGNE 3 =================
\begin{minipage}[t]{0.58\textwidth}
\begin{bloc}{\num{6}Vocabulary at a glance}
{\scriptsize
\renewcommand{\arraystretch}{1.15}
\begin{tabular}{@{}L{2.5cm}@{}l@{\ }L{6.2cm}@{}}
$\models$ (entailment) & \fl & Vrai dans tous les modèles : $e$ découle des clauses\\
Definite clause & \fl & \texttt{tête :- corps.} avec une seule tête positive\\
Fact & \fl & Clause sans corps : \texttt{has\_car(t800,t800\_c1).}\\
Unification & \fl & Rendre deux atomes identiques par une substitution\\
Refutation & \fl & Suite de pas SLD qui mène au but vide\\
Selection rule & \fl & Quel littéral on résout (ici : le plus à gauche)\\
Soundness / Completeness & \fl & Correction / complétude de SLD\\
Allowed constant & \fl & Locale à $i$ ou écrite dans $\mathcal{R}$, $H$\\
\end{tabular}}
\end{bloc}
\end{minipage}\hfill
\begin{minipage}[t]{0.4\textwidth}
\begin{bloc}{\num{7}Pitfalls (les pièges)}
\begin{itemize}
  \item \textbf{Sans directions} : \texttt{f(A) :- has\_car(B,C), short(C).} prend n'importe quel wagon, y compris d'un autre client $\to$ le centralisé prouve, le client non.
  \item \textbf{Constante partagée} : si deux clients ont la même constante, un fait de l'un devient atteignable depuis l'autre.
  \item \textbf{Règle partagée avec constante locale} : casse l'argument ``allowed''.
  \item \textbf{Mauvais ordre} : \texttt{short(B), has\_car(A,B)} appelle \texttt{short} avec \texttt{B} libre.
\end{itemize}
\end{bloc}
\end{minipage}

\vspace{5pt}

\begin{tcolorbox}[colback=marine, colframe=marine, arc=0pt, boxrule=0pt,
  halign=center, top=3pt, bottom=3pt]
\color{white}\small A proof = clear assumptions + a tool (SLD) + an invariant checked at every step.
\end{tcolorbox}

% =====================================================================
%                    PAGES DÉTAILLÉES
% =====================================================================
\detailpages

{\fontsize{24}{26}\selectfont\bfseries\textcolor{marine}{COMPRENDRE}~LA PREUVE}\\[2pt]
{\normalsize Du plus simple au plus technique. Exemple fil rouge : \texttt{t800} (positif, client~2)
et \texttt{t184} (négatif, client~3) du dataset Trains. Les termes de théorie sont donnés en anglais
(\textcolor{marine}{\textbf{gras}}), avec leur traduction entre parenthèses.}

% ---------------------------------------------------------------- 1
\titrebox{1}{Logical building blocks (les briques logiques)}

\st{Definite clause, fact, goal}
\begin{itemize}
  \item \EN{Fact}{fait} : une clause sans corps. \texttt{has\_car(t800,t800\_c1).}
  \item \EN{Rule}{règle} : \texttt{tête :- corps.} Se lit ``si le corps est vrai, alors la tête l'est''.
        \texttt{f(A) :- has\_car(A,B), short(B).} : un train $A$ est positif s'il a un wagon $B$ court.
  \item \EN{Definite clause}{clause définie} : une règle avec une seule tête positive (pas de négation).
        Les hypothèses de Popper dans la thèse sont de ce type.
  \item \EN{Goal}{but} : ce qu'on veut prouver. \texttt{?- f(t800).} Écrit $\leftarrow f(\texttt{t800})$ en logique.
\end{itemize}

\st{Entailment $\models$ and monotonicity}
$\mathcal{B}\cup H\models e$ veut dire : $e$ est vrai dans \textbf{tous} les modèles de $\mathcal{B}\cup H$.
La logique des clauses définies est \EN{monotone}{monotone} : si $A\models e$ et $A\subseteq A'$, alors $A'\models e$.
\ce{c'est exactement le sens facile du lemme : le client a $\mathcal{B}_i\subseteq\mathcal{B}$.}

% ---------------------------------------------------------------- 2
\titrebox{2}{Unification (l'unification)}
\textbf{Unifier} deux atomes, c'est trouver une \EN{substitution}{substitution} $\theta$ qui les rend identiques.
\begin{center}
\begin{tabular}{@{}l l@{}}
atome du but & \texttt{has\_car(t800, B)}\\
fait du BK & \texttt{has\_car(t800, t800\_c4)}\\
substitution & $\theta=\{\texttt{B}/\texttt{t800\_c4}\}$ \quad (après application, les deux atomes sont identiques)\\
\end{tabular}
\end{center}
Contre-exemple : \texttt{has\_car(t800,B)} et \texttt{has\_car(t184,t184\_c1)} ne s'unifient pas, car
\texttt{t800} $\neq$ \texttt{t184}.
\ce{c'est l'unification qui empêche un but sur \texttt{t800} d'utiliser un fait sur \texttt{t184} \emph{tant que l'argument d'entrée est lié}.}

% ---------------------------------------------------------------- 3
\titrebox{3}{SLD resolution : la méthode de preuve de Prolog}
La résolution SLD est la procédure qu'utilise Prolog pour prouver un but. On part du but que l'on veut
prouver, on cherche une clause dont la tête s'unifie avec lui, et on remplace le but par le corps de la
clause, que l'on prouve à son tour, ainsi de suite.\par\vspace{3pt}
(\textbf{SLD} = \textbf{S}elective \textbf{L}inear resolution for \textbf{D}efinite clauses.)

\st{Sur l'exemple}
\begin{itemize}
  \item but : \texttt{f(t800)}
  \item on utilise la clause \texttt{f(A) :- has\_car(A,B), short(B)} avec \texttt{A = t800}, et le but devient
        \texttt{has\_car(t800,B), short(B)}
  \item on résout \texttt{has\_car(t800,B)} avec un fait, par exemple \texttt{has\_car(t800,t800\_c4)}, et il reste
        \texttt{short(t800\_c4)}
  \item on résout \texttt{short(t800\_c4)} avec le fait correspondant, et il ne reste plus rien (le \textbf{but vide})
\end{itemize}

\st{La même chose, étape par étape, avec le backtracking}
\begin{tabular}{@{}l L{5cm} L{9cm}@{}}
\textbf{Step} & \textbf{Current goal} & \textbf{What happens}\\\hline
0 & $\leftarrow$ \texttt{f(t800)} & clause de $H$, \texttt{A=t800}\\
1 & $\leftarrow$ \texttt{has\_car(t800,B), short(B)} & on résout \texttt{has\_car}, premier fait possible : \texttt{B=t800\_c1}\\
2 & $\leftarrow$ \texttt{short(t800\_c1)} & aucun fait \texttt{short(t800\_c1)} : échec, \textbf{backtracking}\\
3 & $\leftarrow$ \texttt{short(t800\_c2)}, puis \texttt{c3} & même échec\\
4 & $\leftarrow$ \texttt{short(t800\_c4)} & fait présent : but \textbf{vide}, \textbf{succès}\\
\end{tabular}

\st{Sur \texttt{f(t184)} (négatif, couvert par erreur)}
\texttt{has\_car(t184,B)} donne \texttt{B=t184\_c1}, puis \texttt{short(t184\_c1)} est un fait : succès immédiat.
Le programme couvre ce négatif (un \EN{false positive}{faux positif}).
\ce{dans les deux cas, chaque fait utilisé parle \emph{uniquement} du train testé. C'est ce qu'on va généraliser.}

% ---------------------------------------------------------------- 4
\titrebox{4}{``Refutation of $\leftarrow e$'' (réfutation)}
Prouver $e$ revient à montrer que l'ajout de sa négation, $\leftarrow e$ (``$e$ est faux''), mène à une
contradiction. Une \EN{SLD refutation}{réfutation SLD} est donc une suite de pas de résolution qui part de
$\leftarrow e$ et arrive au but vide. C'est exactement la dérivation de la section précédente.\par\vspace{3pt}
Dire qu'il y a une réfutation de $\leftarrow e$ dans $\mathcal{B}\cup H$ veut dire que Prolog peut prouver $e$
avec ces clauses.
\ce{réfutation = succès de Prolog. ``Il existe une réfutation'' = ``il existe un chemin de preuve qui réussit''.}

% ---------------------------------------------------------------- 5
\titrebox{5}{``By completeness of SLD resolution''}
\st{Completeness (complétude) et soundness (correction)}
\begin{itemize}
  \item \EN{Soundness}{correction} : s'il existe une réfutation SLD de $\leftarrow e$ dans $P$, alors $P\models e$.
  \item \EN{Completeness}{complétude} : si $P\models e$, alors il existe une réfutation SLD de $\leftarrow e$ dans $P$.
        Cela reste vrai pour \textbf{n'importe quelle} selection rule (littéral de gauche, de droite\ldots).
\end{itemize}
La complétude dit : si $\mathcal{B}\cup H\models e$ (c'est vrai logiquement), alors il existe une réfutation
SLD de $\leftarrow e$. On a donc le droit de remplacer la conséquence logique (difficile à manipuler) par une
\textbf{dérivation concrète}, dont on peut inspecter chaque pas. C'est ce qui permet de suivre les faits
utilisés, un par un.

\st{Où on s'en sert dans la preuve}
\begin{tabular}{@{}L{3.2cm}L{13.2cm}@{}}
\textbf{Completeness} & Au début : $\mathcal{B}\cup H\models e$ est une propriété sémantique, difficile à inspecter. La complétude la remplace par une dérivation concrète, dont on peut examiner chaque clause utilisée.\\[3pt]
\textbf{Soundness} & À la fin : on a une réfutation dans $\mathcal{B}_i\cup H$, et la correction nous dit que cela suffit pour conclure $\mathcal{B}_i\cup H\models e$.\\
\end{tabular}
\ce{la complétude sert à \emph{ouvrir} la preuve, la correction à la \emph{fermer}.}

% ---------------------------------------------------------------- 6
\titrebox{6}{``Always selects the leftmost literal''}
À chaque pas, le but contient plusieurs littéraux (par exemple \texttt{has\_car(t800,B), short(B)}), et il faut
choisir lequel on résout en premier. C'est la \EN{selection rule}{règle de sélection}. Ici on choisit toujours
le littéral le plus à gauche, comme Prolog.\par\vspace{3pt}
La complétude de SLD est valable pour \textbf{n'importe quelle} règle de sélection, donc on peut imposer celle
de notre choix sans perdre de réfutation.

\st{Selection rule $\neq$ search rule}
\begin{itemize}
  \item \textbf{Selection rule} : \emph{quel littéral} on résout dans le but. Ici : le plus à gauche.
  \item \textbf{Search rule} : \emph{dans quel ordre on essaie les clauses}. Prolog : profondeur d'abord, avec backtracking.
        Elle peut boucler, mais ce n'est pas un problème pour nous : on a seulement besoin qu'une réfutation \emph{existe}.
\end{itemize}

% ---------------------------------------------------------------- 7
\titrebox{7}{``With clause bodies written in the order required by Assumption 3''}
Le choix ``le plus à gauche'' n'a de sens que si l'ordre dans lequel les littéraux sont écrits est le bon.
L'\EN{Assumption}{hypothèse} 3 (\textbf{Directions}) dit que chaque argument \texttt{in} d'un littéral est soit
une variable \texttt{in} de la tête, soit une variable \texttt{out} d'un littéral précédent. Pour que
``précédent'' ait un sens, on écrit les corps dans cet ordre :
\begin{moi}
\texttt{f(A) :- has\_car(A,B), short(B).}\quad \textit{\% B est produit par has\_car avant d'être lu par short}
\end{moi}
Avec ce choix, le littéral le plus à gauche a toujours ses entrées déjà connues quand on le résout. C'est ce
que la preuve du lemme exploite : quand on sélectionne un littéral, ses arguments \texttt{in} sont des
\EN{allowed constants}{constantes autorisées}.
\begin{piege}
\textbf{Dans le mauvais ordre} (\texttt{short(B), has\_car(A,B)}), \texttt{short(B)} serait appelé avec
\texttt{B} \emph{inconnu} et pourrait s'unifier avec n'importe quel fait \texttt{short(\ldots)}, y compris
ceux des autres clients.
\end{piege}
\ce{``leftmost'' + ``bodies in the order of the directions'' garantissent que chaque littéral est appelé
avec ses entrées déjà liées à l'exemple.}

% ---------------------------------------------------------------- 8
\titrebox{8}{The proof, line by line (la preuve, ligne par ligne)}

\st{Definitions used in the proof}
\textbf{Allowed constant for client $i$} (constante autorisée) : une constante locale à $i$, ou écrite dans
$\mathcal{R}$ ou dans $H$. Par \textbf{Separated Local Constants}, aucune constante autorisée pour $i$ n'est
locale à un autre client $j$.

\st{Invariant}
\begin{dire}
Chaque fois qu'un littéral est sélectionné, ses arguments \texttt{in} sont des constantes autorisées pour $i$,
et chaque valeur qu'il calcule est aussi une constante autorisée pour $i$.
\end{dire}
On le prouve \textbf{par récurrence} (\EN{by induction}{par récurrence}) \textbf{sur la longueur de la dérivation}.

\st{Base case (cas de base)}
Le but est $f(c)$, avec $c$ une constante locale de $i$ (\textbf{Example Locality}). Sur Trains : \texttt{c = t800}.

\st{Inductive step, Case 1 : le littéral est résolu avec un fait}
\begin{itemize}
  \item Les arguments \texttt{in} du littéral sont des constantes autorisées pour $i$ (\EN{inductive hypothesis}{hypothèse de récurrence}).
  \item Le fait choisi a ces constantes à ses positions \texttt{in}.
  \item \textbf{Si ce fait était chez un autre client $j$} : ses constantes seraient locales à $j$, donc pas
        autorisées pour $i$. Contradiction.
  \item Donc le fait est dans $F_i\cup\mathcal{R}=\mathcal{B}_i$. Ses autres constantes sont locales à $i$ ou
        écrites dans $\mathcal{R}$ : autorisées. L'invariant est conservé.
\end{itemize}
Sur Trains : \texttt{has\_car(t800,B)} a \texttt{t800} en entrée, donc le fait est
\texttt{has\_car(t800,t800\_c4)}, présent chez le client 2. Sa sortie \texttt{t800\_c4} est locale au client 2.

\st{Inductive step, Case 2 : le littéral est résolu avec une règle de $\mathcal{R}\cup H$}
\begin{itemize}
  \item Les variables d'entrée de la tête reçoivent des constantes autorisées.
  \item Dans le corps, chaque argument \texttt{in} est : une variable d'entrée de la tête, une constante de la règle
        (autorisée), ou une variable de sortie d'un littéral \emph{précédent} (résolu avant, donc sa valeur est
        autorisée par hypothèse de récurrence).
  \item Comme toute variable de la tête apparaît dans le corps, les valeurs que la tête calcule sont aussi autorisées.
\end{itemize}

\st{Conclusion}
Toutes les clauses utilisées par la réfutation sont dans $\mathcal{B}_i\cup H$. C'est donc une réfutation SLD
de $\leftarrow e$ dans $\mathcal{B}_i\cup H$, et par \textbf{soundness} : $\mathcal{B}_i\cup H\models e$. \hfill $\square$

% ---------------------------------------------------------------- 9
\titrebox{9}{Counter-examples : pourquoi chaque assumption}
\renewcommand{\arraystretch}{1.3}
\begin{tabular}{@{}L{3.4cm}L{6cm}L{7cm}@{}}
\textbf{Assumption removed} & \textbf{Example} & \textbf{What breaks}\\\hline
Directions & \texttt{quelqu\_un\_a\_fievre(X) :- fever(Y).} & \texttt{Y} libre : il s'unifie avec le fait \texttt{fever(bob)} d'un autre client. Le centralisé prouve, le client non.\\
Separated Local Constants & Deux clients ont tous deux un wagon \texttt{c1} & Un fait du client $j$ peut être atteint depuis un but du client $i$ par cette constante commune.\\
Example Locality & Un exemple dont la constante n'est locale à personne & La dérivation démarre hors de $\mathcal{B}_i$.\\
\end{tabular}

% ---------------------------------------------------------------- 10
\titrebox{10}{Questions qu'un relecteur peut poser (reviewer questions)}

\st{``Pourquoi ne pas simplement supposer la knowledge locality ?''}
\begin{dire}
Parce qu'elle découle de conditions plus simples, qu'on peut vérifier sur les données et le biais :
les constantes sont séparées et les clauses respectent les directions. La démontrer permet de dire
\emph{quand} elle est vraie.
\end{dire}
\begin{direEN}
Because it follows from simpler conditions that can be checked on the data and the bias: constants are
separated and clauses respect the directions. Proving it tells us \emph{when} it holds.
\end{direEN}

\st{``Pourquoi la complétude de SLD, et pas juste le $\models$ ?''}
\begin{dire}
Le $\models$ est une propriété sémantique. La complétude nous donne une dérivation concrète, dont on
peut inspecter chaque clause utilisée. C'est ce qui permet de montrer que tous les faits utilisés sont locaux.
\end{dire}
\begin{direEN}
Entailment is a semantic property. Completeness gives us a concrete derivation whose every used clause we can
inspect, which lets us show that all facts used are local.
\end{direEN}

\st{``Que se passe-t-il si $\mathcal{R}$ contient des règles ?''}
\begin{dire}
La preuve les traite au Cas 2 : tant qu'elles respectent les directions et ne contiennent pas de constante locale,
les valeurs qu'elles calculent restent autorisées.
\end{dire}
\begin{direEN}
The proof handles them in Case 2: as long as they respect the directions and contain no local constant, the values
they compute remain allowed.
\end{direEN}

\st{``Est-ce que cela dépend du solveur ou de Prolog ?''}
\begin{dire}
Non : la preuve porte sur l'existence d'une réfutation, pas sur la stratégie de recherche de Prolog.
\end{dire}
\begin{direEN}
No: the proof is about the existence of a refutation, not about Prolog's search strategy.
\end{direEN}

% ---------------------------------------------------------------- 11
\titrebox{11}{Proof vocabulary (le vocabulaire des preuves)}
\renewcommand{\arraystretch}{1.25}
\begin{tabular}{@{}L{3.6cm}L{3.4cm}L{9.4cm}@{}}
\textbf{English} & \textbf{Français} & \textbf{Meaning}\\\hline
Definition & définition & fixe un objet ou une notation, sans preuve\\
Assumption & hypothèse (de travail) & condition qu'on suppose vraie ; sans elle le résultat peut être faux\\
Lemma & lemme & résultat intermédiaire, une brique pour un résultat plus important\\
Proposition & proposition & résultat intéressant en lui-même, moins central qu'un théorème\\
Theorem & théorème & résultat principal\\
Corollary & corollaire & conséquence directe d'un théorème\\
Remark & remarque & commentaire, sans nouveau résultat\\
Proof & preuve, démonstration & enchaînement logique qui établit un énoncé\\
Invariant & invariant & propriété vraie à chaque étape d'un processus\\
Induction & récurrence & preuve en deux temps : cas de base, puis pas inductif\\
Base case / Inductive step & cas de base / pas de récurrence & les deux parties d'une preuve par récurrence\\
Inductive hypothesis & hypothèse de récurrence & ce qu'on suppose vrai pour l'étape $n$ pour montrer $n+1$\\
Contradiction & contradiction, absurde & supposer le contraire et aboutir à une impossibilité\\
If and only if ($\iff$) & si et seulement si & équivalence : les deux sens\\
Entailment ($\models$) & conséquence logique & vrai dans tous les modèles\\
Soundness / Completeness & correction / complétude & réfutation $\Rightarrow$ conséquence / conséquence $\Rightarrow$ réfutation\\
\end{tabular}

\vspace{8pt}
\begin{tcolorbox}[colback=marine, colframe=marine, arc=0pt, boxrule=0pt,
  halign=center, top=4pt, bottom=4pt]
\color{white}\small Si je sais refaire la trace de \texttt{t800} et dire pourquoi chaque assumption est nécessaire, je sais la preuve.
\end{tcolorbox}

\end{document}
